%
% kehrwert.tex -- Graph der Funktion f(x) = 1/x
%
% (c) 2026 Damien Flury, OST Ostschweizer Fachhochschule
%
\documentclass[tikz]{standalone}
\usepackage{amsmath}
\usepackage{times}
\usepackage{txfonts}
\usepackage{pgfplots}
\usetikzlibrary{arrows,math}
\definecolor{darkred}{rgb}{0.8,0,0}
\begin{document}
\def\skala{1}
\begin{tikzpicture}[>=latex,thick,scale=\skala]

% Achsen
\draw[->] (-4.3,0) -- (4.5,0) coordinate[label={$x$}];
\draw[->] (0,-3.3) -- (0,3.5) coordinate[label={right:$y$}];

% Achsenbeschriftung
\foreach \x in {-4,-3,-2,-1,1,2,3,4}{
	\draw (\x,-0.05) -- (\x,0.05);
	\node at (\x,0) [below] {$\x$};
}
\foreach \y in {-3,-2,-1,1,2,3}{
	\draw (-0.05,\y) -- (0.05,\y);
	\node at (0,\y) [left] {$\y$};
}

% Flächen unter dem Integranden, ungleichmässig um die Polstelle gewählt
\fill[darkred!30,opacity=0.5]
	(-3.5,0) -- plot[domain=-3.5:-0.6,samples=100] ({\x},{1/\x}) -- (-0.6,0) -- cycle;
\fill[darkred!30,opacity=0.5]
	(0.4,0) -- plot[domain=0.4:2.5,samples=100] ({\x},{1/\x}) -- (2.5,0) -- cycle;

% negativer Ast
\draw[color=darkred,line width=1.4pt]
	plot[domain=-4.2:-0.303,samples=200] ({\x},{1/\x});
% positiver Ast
\draw[color=darkred,line width=1.4pt]
	plot[domain=0.303:4.2,samples=200] ({\x},{1/\x});

\node[color=darkred] at (2.3,1.1) {$f(x)=\dfrac{1}{x}$};

\end{tikzpicture}
\end{document}
